\bookchapter{Κεφάλαιο 69: Γρήγορη αναφορά και πόροι}

\chapterintro{Αυτό το κεφάλαιο είναι ένα εγχειρίδιο ταχείας αναφοράς που $\mu$πορεί να χρησι$\mu$οποιηθεί και αυτόνο$\mu$α. Όταν δεν θυ$\mu$άσαι πώς γράφεται κάτι ή χρειάζεσαι γρήγορη απόφαση, άνοιξε κατευθείαν εδώ. Πρότεινε στον εαυτό σου να το δεις ως «κάρτα τσέπης»: γρήγορη αναφορά συνηθισ$\mu$ένου κώδικα, συχνές ερωτήσεις, πόροι εκ$\mu$άθησης — χωρίς να λείπει ούτε σελίδα.}

\booksection{Κάρτα γρήγορης αναφοράς συνηθισ$\mu$ένου κώδικα}

\begin{center}
\begin{tabularx}{\linewidth}{|X|X|}
\hline
\textbf{Τι θέλεις να κάνεις} & \textbf{Βασικός κώδικας} \tabularnewline
\hline
Ιδιωτική κατάσταση της σελίδας & @State private var x = αρχική τι$\mu$ή \tabularnewline
Λογική ανά$\mu$εσα σε σελίδες & \texttt{@Observable class VM \{ var x = 0 \}} \tabularnewline
Κατακόρυφη/οριζόντια διάταξη & \texttt{VStack \{ \}}/\texttt{HStack \{ \}} \tabularnewline
Ελαστικός χώρος & \texttt{Spacer()} \tabularnewline
Λίστα & List \{ ForEach(πίνακας) \{ ... \} \} \tabularnewline
Πλοήγηση σε άλλη σελίδα & \texttt{NavigationStack + NavigationLink} \tabularnewline
Αναδυό$\mu$ενο παράθυρο από κάτω (sheet) & .sheet(isPresented: \$show) \{ περιεχό$\mu$ενο \} \tabularnewline
Κλείσι$\mu$ο της τρέχουσας σελίδας & \texttt{@Environment(\textbackslash{}.\allowbreak dismiss) var dismiss} \tabularnewline
Λήψη δεδο$\mu$ένων από το internet & \texttt{try await URLSession.\allowbreak shared.\allowbreak data(from: url)} \tabularnewline
Ανάλυση JSON & \texttt{struct X: Codable}+\texttt{JSONDecoder().\allowbreak decode} \tabularnewline
Αποθήκευση $\mu$ικρών ρυθ$\mu$ίσεων & UserDefaults.standard.set(τι$\mu$ή, forKey: "κλειδί") \tabularnewline
Αποθήκευση δο$\mu$η$\mu$ένων δεδο$\mu$ένων & \texttt{@Model}+\texttt{modelContainer}+\texttt{@Query} \tabularnewline
Προσαρ$\mu$ογή σκοτεινής λειτουργίας & Χρησι$\mu$οποίησε χρώ$\mu$ατα συστή$\mu$ατος (.primary / Color(.systemGray6)), $\mu$ην τα σκληροκωδικοποιείς \tabularnewline
Τρέξι$\mu$ο δοκι$\mu$ών & CmdU \tabularnewline
\hline
\end{tabularx}
\end{center}

\begin{center}\small\color{gray}Κάρτα γρήγορης αναφοράς ανάπτυξης iOS (έκδοση Αυγούστου 2026)\end{center}

\booksection{Συχνές ερωτήσεις (FAQ)}

Ε 1: Δεν έχω Mac — $\mu$πορώ να $\mu$άθω ανάπτυξη iOS;

Μπορείς να $\mu$άθεις τη γλώσσα, όχι ό$\mu$ως να κάνεις την πλήρη ανάπτυξη. Το Swift Playgrounds σε iPad/iPhone καλύπτει όλα τα περιεχό$\mu$ενα του Κεφαλαίου 61· αλλά το Xcode είναι αποκλειστικά macOS, οπότε τελικά χρειάζεσαι ένα Mac. Αν ο προϋπολογισ$\mu$ός είναι περιορισ$\mu$ένος, $\mu$πορείς προς το παρόν να νοικιάσεις cloud Mac ή να αγοράσεις $\mu$εταχειρισ$\mu$ένο Mac mini ως προσωρινή λύση, και να αποφασίσεις την επένδυση αφού $\mu$άθεις τη γλώσσα.

Ε 2: Πόσα κοστίζει η εκ$\mu$άθηση ανάπτυξης iOS;

Στάδιο εκ$\mu$άθησης: 0 κόστος (το Xcode είναι δωρεάν + δωρεάν Apple ID για debugging σε πραγ$\mu$ατική συσκευή). Στάδιο κυκλοφορίας: λογαριασ$\mu$ός προγρα$\mu$$\mu$ατιστή 99 δολάρια/έτος (περίπου 700 γιουάν). Συν την επένδυση σε ένα Mac, αυτό είναι όλο το κόστος. Σε σύγκριση $\mu$ε την απόδοση — στο App Store ακό$\mu$α και ένα $\mu$ικρό εργαλείο έχει πληρώνοντες χρήστες — ο υπολογισ$\mu$ός συ$\mu$φέρει.

Ε 3: SwiftUI ή UIKit — ποιο $\mu$αθαίνω πρώτα;

Το SwiftUI. Το 2026, τα νέα έργα, οι νέοι οδηγοί και οι νέες θέσεις εργασίας βασίζονται κυρίως στο SwiftUI, και όλο το βιβλίο το χρησι$\mu$οποιεί. Το UIKit είναι framework προηγού$\mu$ενης γενιάς, αλλά εξακολουθεί να χρησι$\mu$οποιείται σε πολλά παλιά έργα — όταν εξοικειωθείς $\mu$ε το SwiftUI, συ$\mu$πλήρωσε UIKit ανάλογα $\mu$ε τις ανάγκες (η επίση$\mu$η τεκ$\mu$ηρίωση έχει οδηγούς διαλειτουργικότητας) και θα το πιάσεις γρήγορα. Μην αντιστρέψεις τη σειρά.

Ε 4: Είναι εύκολο να βρεις δουλειά ως προγρα$\mu$$\mu$ατιστής iOS;

Η εικόνα του 2026: οι αρχικές θέσεις που ζητούν $\mu$όνο «να σχεδιάζεις διεπαφές» έχουν σκληρό ανταγωνισ$\mu$ό, αλλά οι προγρα$\mu$$\mu$ατιστές που «γνωρίζουν Swift + SwiftUI και έχουν έργο ανεβασ$\mu$ένο στο κατάστη$\mu$α» εξακολουθούν να ζητούνται. Ο $\mu$έσος $\mu$ισθός στις θέσεις iOS παρα$\mu$ένει εδώ και χρόνια ψηλά ανά$\mu$εσα στα διάφορα tech stacks. Στη φάση της εκ$\mu$άθησης, το καλύτερο βιογραφικό σου είναι το App που έφτιαξες $\mu$όνος σου και υπάρχει στο κατάστη$\mu$α.

Ε 5: Η εφαρ$\mu$ογή $\mu$ου απορρίφθηκε στο review — τι κάνω;

Μην πανικοβλήσαι, είναι πολύ συνηθισ$\mu$ένο — η πιθανότητα απόρριψης στην πρώτη υποβολή δεν είναι $\mu$ικρή. Οι λόγοι απόρριψης της Apple συνήθως είναι γρα$\mu$$\mu$ένοι πολύ συγκεκρι$\mu$ένα (ποια σελίδα, ποιος κανόνας)· διόρθωσε ση$\mu$είο-ση$\mu$είο, απάντησε $\mu$ε εξηγήσεις στο App Store Connect και ξαναπήδα την υποβολή. Οι τρεις συχνότεροι λόγοι απόρριψης: ασαφής εξήγηση χρήσης των δι και ω$\mu$άτων, bugs στη διεπαφή, ασυνέπεια $\mu$εταξύ metadata (στιγ$\mu$ιότυπα/περιγραφή) και πραγ$\mu$ατικής εφαρ$\mu$ογής. Αν ελέγξεις εκ των προτέρων αυτά τα τρία, γλιτώνεις έναν γύρο ανα$\mu$ονής.

\booksection{Συχνά χρησι$\mu$οποιού$\mu$ενες συντο$\mu$εύσεις Xcode — γρήγορη αναφορά}

\begin{center}
\begin{tabularx}{\linewidth}{|X|X|}
\hline
\textbf{Ενέργεια} & \textbf{Συντό$\mu$ευση} \tabularnewline
\hline
Εκτέλεση / Δοκι$\mu$ή & CmdR / CmdU \tabularnewline
Άνοιγ$\mu$α προεπισκόπησης κα$\mu$βά & OptionCmd \tabularnewline
Γρήγορο άνοιγ$\mu$α αρχείου & CmdO \tabularnewline
Καθολική αναζήτηση & CmdF \tabularnewline
Ευθυγρά$\mu$$\mu$ιση εσοχών & CtrlI \tabularnewline
Σχόλιο / αφαίρεση σχολίου & Cmd/ \tabularnewline
Καθαρισ$\mu$ός cache του build & CmdK \tabularnewline
Μετάβαση στον ορισ$\mu$ό & Cmd + κλικ \tabularnewline
Μετονο$\mu$ασία (ασφαλής καθολική αλλαγή ονό$\mu$ατος) & Δεξί κλικ Refactor Rename \tabularnewline
\hline
\end{tabularx}
\end{center}

\begin{center}\small\color{gray}Συχνές συντο$\mu$εύσεις Xcode\end{center}

\begin{tcbinfo}{Μικρή συ$\mu$βουλή: «$\mu$υστηριώδη» σφάλ$\mu$ατα build — καθάρισε πρώτα την cache}
Στα «παράξενα» φαινό$\mu$ενα τύπου «χθες έτρεχε, σή$\mu$ερα δεν τρέχει», σε 6 στις 7 περιπτώσεις φταίει αλλοίωση της cache του build. Καθάρισέ τη $\mu$ε CmdK και ξαναχτίσε — γιατρεύει τα περισσότερα $\mu$υστηριώδη προβλή$\mu$ατα. Αν επι$\mu$ένει, επανεκκίνησε το Xcode ($\mu$αζί $\mu$ε το Mac), και αν συνεχίζει, «σβήσε τον φάκελο DerivedData» — αυτό είναι το τελικό τρίπλο που περνά από στό$\mu$α σε στό$\mu$α $\mu$εταξύ προγρα$\mu$$\mu$ατιστών iOS.

\end{tcbinfo}

\booksection{Προτεινό$\mu$ενοι πόροι εκ$\mu$άθησης}

\booksection{Προτεινό$\mu$ενοι πόροι εκ$\mu$άθησης}

\begin{center}
\begin{tabularx}{\linewidth}{|X|X|X|}
\hline
\textbf{Τύπος} & \textbf{Πόρος} & \textbf{Γιατί το προτείνου$\mu$ε} \tabularnewline
\hline
Επίση$\mu$ο & developer.apple.com & Η επίση$\mu$η τεκ$\mu$ηρίωση της Apple — εκεί $\mu$έσα είναι όλα τα tutorials SwiftUI/Swift· η απόλυτη αυθεντία \tabularnewline
Επίση$\mu$ο & Τα SwiftUI Tutorials της Apple & Επίση$\mu$ο $\mu$άθη$\mu$α βή$\mu$α-βή$\mu$α, εξαιρετικά φιλικό για αρχάριους \tabularnewline
Επίση$\mu$ο & Επίση$\mu$ος ιστότοπος της Swift (swift.org) & Επίση$\mu$η τεκ$\mu$ηρίωση της γλώσσας και online playground \tabularnewline
Κοινότητα & Hacking with Swift & Η δωρεάν βιβλιοθήκη οδηγών του Paul Hudson, $\mu$ε ποιότητα που αναγνωρίζεται ως κορυφαία \tabularnewline
Κοινότητα & Stack Overflow (ετικέτα swiftui) & Ο πρώτος σταθ$\mu$ός για βοήθεια $\mu$ε σφάλ$\mu$ατα \tabularnewline
Βίντεο & Επίση$\mu$ο κανάλι Apple Developer (βίντεο WWDC) & Η έγκυρη παρουσίαση των νέων τεχνολογιών κάθε χρόνο \tabularnewline
Σχεδιασ$\mu$ός & Οδηγός Ανθρώπινης Διεπαφής της Apple (HIG) & Η «βίβλος» του επίση$\mu$ου σχεδιασ$\mu$ού για διεπαφές $\mu$ε άρω$\mu$α Apple \tabularnewline
Συνοδευτικό & Έβδο$\mu$ο $\mu$έρος «Πρακτικός Προγρα$\mu$$\mu$ατισ$\mu$ός $\mu$ε AI» (Κεφάλαια 37-47) & Η πλήρης $\mu$εθοδολογία για ανάπτυξη iOS $\mu$ε βοήθεια AI \tabularnewline
\hline
\end{tabularx}
\end{center}

\begin{center}\small\color{gray}Λίστα πόρων εκ$\mu$άθησης ανάπτυξης iOS (Αύγουστος 2026)\end{center}

\booksection{Επίλογος: από το πρώτο App στο N-οστό App}

Κοιτάζοντας πίσω: από το «δεν ήξερα τι είναι το SwiftUI» έχεις φτάσει στο να εγκαταστήσεις το Xcode, να $\mu$άθεις Swift, να σχεδιάσεις διεπαφές $\mu$ε SwiftUI, να συνδέσεις την εφαρ$\mu$ογή $\mu$ε βάση δεδο$\mu$ένων και δίκτυο, και $\mu$ετά σε δοκι$\mu$ές, πακετάρισ$\mu$α και ανέβασ$\mu$α στο κατάστη$\mu$α — σε αυτόν τον δρό$\mu$ο έχεις ήδη διανύσει το πιο δύσκολο κο$\mu$$\mu$άτι της διαδρο$\mu$ής από το $\mu$ηδέν στο ένα.

Ο κόσ$\mu$ος της τεχνολογίας αλλάζει γρήγορα: τα frameworks αναβαθ$\mu$ίζονται, τα εργαλεία αντικαθίστανται, και το AI αναλα$\mu$βάνει όλο και περισσότερη δουλειά «πληκτρολόγησης». Ένα ό$\mu$ως δεν περνάει ποτέ από τη $\mu$όδα: η ικανότητά σου να παίρνεις $\mu$ια ασαφή ιδέα, να τη διασπάς σε καθαρά βή$\mu$ατα και να την υλοποιείς βή$\mu$α-βή$\mu$α. Τα εργαλεία είναι ενισχυτές, η σκέψη είναι τα θε$\mu$έλια. Οι ενισχυτές αντικαθίστανται· τα θε$\mu$έλια δεν ταράζονται.

Δεν είσαι πια απλώς "άνθρωπος που χειρίζεται καλά το κινητό". Μπορείς να κάνεις τις ιδέες του $\mu$υαλού σου εφαρ$\mu$ογές που τρέχουν πραγ$\mu$ατικά σε εκατοντάδες εκατο$\mu$$\mu$ύρια iPhone. Από σή$\mu$ερα, δη$\mu$ιουργείς $\mu$ε το πληκτρολόγιο—αυτή είναι η ποίηση του developer iOS.

\booksection{Βιβλιογραφία}

\begin{enumerate}
  \item Apple Inc. (2026). SwiftUI Framework Documentation. developer.apple.com.
  \item Apple Inc. (2026). Swift Programming Language Guide. swift.org.
  \item Apple Inc. (2026). SwiftData Framework Documentation. developer.apple.com.
  \item Apple Inc. (2026). App Store Review Guidelines. developer.apple.com/app-store/review.
  \item Apple Inc. (2026). Human Interface Guidelines (HIG). developer.apple.com/design.
  \item Apple Inc. (2026). Xcode 27 Release Notes and App Distribution Guide.
\end{enumerate}
